\section{Degree Upper Bound}
\label{sec:upper-bound}

In this section we give a simple upper bound on the SoS refutation
degree for $\circuit_s(f)$.  Specifically, for functions $f$ that have
no circuit of size $s$, we show that there is an SoS refutation of
$\circuit_s(f)$ of degree $O(s)$, essentially matching our
$\Omega(s^{1-\epsilon})$ lower bound.

At a high level, the logic behind the refutation is as follows: first,
we show that SoS in degree $O(s)$ can derive that the $\Theta(s^2)$
structure variables of $\circuit_s(f)$ must uniquely describe a
circuit of size $s$.  Then, we also show that for any fixed circuit,
SoS can in degree $O(s)$ derive that the circuit does not compute $f$
by ``evaluating'' the circuit on some (non-deterministically chosen)
input where the output differs from $f$.

To make this precise, we first define a set of monomials, which correspond
to circuits of size $s$. A multilinear monomial $m$ is a \emph{circuit
  monomial} if for every gate $v \in [s]$ it holds that
\begin{enumerate}
\item exactly one of the variables $\isneg(v), \isor(v)$ or $\isand(v)$ occurs
  in $m$,
\item for $a \in \set{1,2}$ exactly one of the variables
  $\isfromconst(v, a)$, $\isfromvar(v, a)$ or $\isfromgate(v, a)$ occurs
  in $m$,
\item for $a \in \set{1,2}$ exactly one of the variables
  $\constval(v, a)$ or $\overline{\constval(v, a)}$ occurs in $m$,
\item for $a \in \set{1,2}$ exactly one of the variables
  $\set{\isvar(v,a,i) \mid i \in [n]}$
  occurs in $m$, and
\item for $a \in \set{1,2}$ and $v > 1$, exactly one of the variables
  $\set{\isgate(v,a,u) \mid u < v}$
  occurs in $m$, and
  \item no other variables occur in $m$ than the ones described
    above.
\end{enumerate}
We denote by $\calM_s$ the set of circuit monomials. We first show
that SoS can derive in degree $O(s)$ the polynomial
$\sum_{m \in \calM_s} m - 1$; this corresponds to SoS proving that
the structure variables uniquely describe a circuit and does not use
anything about $f$.  Then in a second step we show that for every
$m \in \calM_s$, SoS can derive $-m$ in degree $O(s)$; this
corresponds to SoS proving that the circuit described by $m$ does not
compute $f$ correctly.  Summing these two parts up yields an SoS
derivation of $-1$, i.e., a refutation of the $\circuit_s(f)$ formula.

\pa{Here we still talk about SoS  derivations so we can't remove it from prelims}

\paragraph{Deriving~$\bm{\sum_{m \in \calM_s} m - 1}$.} We proceed
by induction on $s$. Note that $\calM_0 = \set{1}$ and hence the base
case is trivial. Suppose we have an SoS derivation of
$\sum_{m \in \calM_s} m - 1$. For every monomial $m \in \calM_s$ we
add the polynomial
\begin{align}
  m \cdot \big(\isneg(v) +  \isor(v) + \isand(v) - 1\big)
\end{align}
to the derivation (note that the second term is \cref{eq:ax-fn}). This
gives us an SoS derivation of $\sum_{m \in \calM'_s} m - 1$, where
$$\calM'_s = \bigcup_{m \in \calM_s} \set{m \cdot \isneg(v), m \cdot \isor(v), m \cdot \isand(v)}\eqperiod$$ We can continue in the same manner with
\cref{eq:ax-from} and \cref{eq:ax-in-gt} to finally obtain an SoS
derivation of $\sum_{m \in \calM_{s+1}} m - 1$. Clearly this
derivation requires degree at most $O(s)$, as for each gate there are
at most $7$ variables in every monomial from $\calM_s$.

\paragraph{Deriving $-m$ for $m \in \calM_s$.}
Let $C$ be the circuit that corresponds to the monomial $m$ and let
$\alpha \in \set{0,1}^n$ be such that $f(\alpha) \neq C(\alpha)$ (by
the assumption that $f$ does not have a circuit of size $s$, such an
$\alpha$ exists).  Suppose $C(\alpha) = b$ but $f(\alpha) = 1-b$.

We construct a degree $O(s)$ SoS proof of the fact that
$C(\alpha) = b$. That is, we are going to show that the polynomial
$p_s = m \cdot \left( \outwire_\alpha(s) - b\right)$ can be written as
\begin{align}
  \label{eq:-m-derivation}
  p_s = \sum_{i=1}^t r_i \cdot q_i \eqcomma
\end{align}
for some parameter $t$, axioms $q_1, \ldots, q_t$ and some polynomials
$r_1, \ldots, r_t$, such that $\deg(r_i \cdot q_i) = O(s)$ for all
$i$.  Note that given this, we can then easily derive $-m$ by
subtracting the polynomial
$m \cdot \left( \outwire_\alpha(s) - f(\alpha)\right) = m \cdot \left(
  \outwire_\alpha(s) - 1 + b\right)$ (this is $m$ multiplied by
\cref{eq:ax-correct}), yielding a derivation of $m \cdot (1 - 2b)$
which is either $m$ or $-m$ depending on $b$; in the former case it
can be multiplied by $-1$ to yield $-m$.  Note that here it is
important that the derivation \eqref{eq:-m-derivation} is a
Nullstellensatz derivation, not using any Sum-of-Squares part, since
otherwise it would not be possible to multiply it by $-1$.

Let us thus see how to derive $m \cdot \left( \outwire_\alpha(s) - b\right)$.
We do this by structural induction over the circuit: for every gate
$v$ we are going to construct an SoS proof of the fact that the
circuit rooted at $v$ outputs the bit $b_v$ on input $\alpha$. In
other words, an SoS derivation of the polynomial
$m \cdot \left(\outwire_\alpha(v) - b_v\right)$.


Let us explain how to construct an SoS proof $p_v$. Consider a gate
$v$ in the circuit. Depending on the function computed at $v$ and how
the wires of $v$ are connected we construct $p_v$ slightly differently.
As a first step let us construct SoS proofs $q_1$ and $q_2$ of the
fact that on input $\alpha$ the bit $c_a$, $a \in \set{1,2}$, is
carried on wire $a$ to the gate $v$. That is, the polynomial $q_a$
should simplify to $m \cdot \big(\inwire_\alpha(v, a) - c_a\big)$. In
the following we explain how to precisely define $q_a$ depending on
what the wire is connected to. Note that not a lot is going on -- we
are mostly just multilinearizing using the Boolean axioms.

If $m$ is of the form $m = m' \cdot \isfromconst(v,a) \cdot \constval(v, a)$ for some monomial $m'$, i.e., wire $a$
is connected to the constant $c_a = 1$ in the circuit described by $m$, then we
can derive $q_a = m \cdot \big(\inwire_{\alpha}(v,a) - 1\big)$
by the identity
\begin{align}
  \nonumber
  q_a &= m' \cdot \constval(v,a) \cdot \isfromconst(v,a) \cdot \big(\inwire_{\alpha}(v,a) -
        \constval(v, a)\big) +\\
      &\qquad m' \cdot \isfromconst(v,a) \cdot
        \big( \constval(v,a)^2 - \constval(v,a)\big) \eqcomma
\end{align}
a linear combination of \cref{eq:inwire-const} and the Boolean axiom on $\constval(v,a)$.  Similarly if $m = m' \cdot \isfromconst(v,a) \cdot \overline{\constval(v,a)}$ (i.e., $c_a = 0$) we can derive $q_a = m \cdot \inwire_{\alpha}(v,a)$ by
\begin{align}
  \nonumber
  q_a &= m' \cdot \overline{\constval(v,a)} \cdot \isfromconst(v,a) \cdot \big(\inwire_{\alpha}(v,a) -
        \constval(v, a)\big) +\\\nonumber
      &\qquad m' \cdot \isfromconst(v,a) \cdot \constval(v,a) \cdot
         \\\nonumber
      &\qquad\qquad\qquad \big( 1 - \constval(v,a) - \overline{\constval(v,a)}\big) +\\
      &\qquad m' \cdot \isfromconst(v,a) \cdot (\constval(v,a)^2 - \constval(v,a))
        \eqcomma
\end{align}
where we additionally use the negation axiom on $\constval(v,a)$.

Next, if $a$ is connected to an input $i$ and $m$ is of the form
$m = m' \cdot \isfromvar(v, a) \cdot \isvar(v, a, i)$ (so that $c_a = \alpha_i$), then
\begin{align}
  q_a = m \cdot (\inwire_{\alpha}(v,a) - \alpha_i) &= m' \cdot \isfromvar(v, a) \cdot \isvar(v, a, i) \cdot
          \big( \inwire_\alpha(v, a) - \alpha_i\big) \eqcomma
\end{align}
which is a multiple of \cref{eq:inwire-var}.  Lastly, if
$a$ is connected to a gate $u$ and $m$ is of the form
$m = m' \cdot \isfromgate(v, a) \cdot \isgate(v, a, u)$ (i.e., $c_a = b_u$), then
\begin{align}
  \nonumber
  q_a = m \cdot(\inwire_\alpha(v,a) - b_u) &= m' \cdot
          \isfromgate(v, a) \cdot \isgate(v, a, u) \cdot
          \big(\inwire_\alpha(v, a) - \outwire_\alpha(u)\big)\\
                                           &\qquad+ m \cdot (\outwire_\alpha(u) - b_u)
                                             \eqcomma
\end{align}
where the first term is a multiple of \cref{eq:inwire-gate}, and the
second term is the polynomial $p_u$ which, by induction, we assume has
already been derived in degree $O(s)$.

Given the two SoS proofs $q_1$ and $q_2$ we are ready to construct the
SoS proof $p_v = m \cdot (\outwire_\alpha(v) - b_v)$. As mentioned earlier
we do a case distinction over the funtion computed at $v$.
\begin{enumerate}
\item \emph{$v$ is a \emph{not} gate ($m = m' \cdot \isneg(v)$ and $b_v = 1-c_1$).} We have the derivation
  \begin{align}
    \nonumber
    p_v &=  m' \cdot \isneg(v) \cdot
             \big(\outwire_{\alpha}(v) - \overline{\inwire_{\alpha}(v, 1)}\big) \\\nonumber
           &\qquad + m \cdot
             \big(\overline{\inwire_{\alpha}(v, 1)} - 1 + \inwire_{\alpha}(v, 1)\big)\\
           &\qquad - q_1\eqcomma
  \end{align}
  where the first line uses \cref{eq:ax-neg} and the
  second line uses the negation axiom for $\inwire_{\alpha}(v, 1)$.
\item \emph{$v$ is an \emph{or} gate ($m = m' \cdot \isor(v)$, and $b_v = 1-(1-c_1)(1-c_2)$.} We have the derivation
  \begin{align}
    \nonumber
    p_v =  & m' \cdot \isor(v) \cdot 
             \Big(
             \outwire_{\alpha}(v) -
             \big(
             1 -
             \overline{\inwire_{\alpha}(v, 1)} \cdot
             \overline{\inwire_{\alpha}(v, 2)}
             \big)
             \Big)\\\nonumber
           & -
             m \cdot
             \overline{\inwire_\alpha(v, 1)} \cdot
             \big(
             \overline{\inwire_{\alpha}(v, 2)}
             - 1 + \inwire_{\alpha}(v, 2)
             \big)\\\nonumber
           & +
             m \cdot
             \big(
             \inwire_\alpha(v, 2) - 1
             \big)
             \cdot
             \big(
             \overline{\inwire_{\alpha}(v, 1)}
             - 1 + \inwire_{\alpha}(v, 1)
             \big)\\\nonumber
           & +
             \big(
             1 - \inwire_\alpha(v, 1)
             \big)
             \cdot
             q_2 \\
           & +
             \big(
             1 - c_2
             \big)
             \cdot
             q_1\eqcomma
  \end{align}
  where the first line uses \cref{eq:ax-or}, and the
  following two lines uses negation axioms.
\item \emph{$v$ is an \emph{and} gate ($m = m' \cdot \isand(v)$, and $b_v = c_1 \cdot c_2$.} We have
  \begin{align}
    \nonumber
    p_v =  & m' \cdot \isand(v) \cdot
             \big(
             \outwire_{\alpha}(v) -
             \inwire_{\alpha}(v, 1) \cdot
             \inwire_{\alpha}(v, 2)
             \big)\\\nonumber
           & +
             \inwire_\alpha(v, 1)
             \cdot
             q_2 \\
           & +
             c_2
             \cdot
             q_1\eqcomma
  \end{align}
  where the first line uses \cref{eq:ax-and}.
\end{enumerate}
This completes the description of the SoS derivation of
$\outwire_\alpha(s) = b$.
%
Observe that the final proof $p_s$ is of degree $O(s)$: in every
inductive step we increase the degree of the proof by at most a
constant.
